\begin{ex}
Chuyển tổng $2 + 2 + 2 + 2$ thành tích $2 \times 4$. Tính kết quả của phép nhân đó.
\shortans{8}
\loigiai{
Ta có: $2 \times 4 = 8$.
}
\end{ex}

\begin{ex}
Chuyển tổng $5 + 5 + 5$ thành tích $5 \times 3$. Tính kết quả của phép nhân đó.
\shortans{15}
\loigiai{
Ta có: $5 \times 3 = 15$.
}
\end{ex}

\begin{ex}
Chuyển tổng $3 + 3 + 3 + 3 + 3$ thành tích. Phép tính có kết quả bằng bao nhiêu?
\shortans{15}
\loigiai{
Ta có: $3 \times 5 = 15$.
}
\end{ex}

\begin{ex}
Chuyển tổng $4 + 4 + 4$ thành tích. Phép tính có kết quả bằng bao nhiêu?
\shortans{12}
\loigiai{
Ta có: $4 \times 3 = 12$.
}
\end{ex}

\begin{ex}
Chuyển tổng $6 + 6$ thành tích. Phép tính có kết quả bằng bao nhiêu?
\shortans{12}
\loigiai{
Ta có: $6 \times 2 = 12$.
}
\end{ex}

\begin{ex}
Số 7 được lấy 3 lần viết thành phép nhân là $7 \times 3$. Tính kết quả.
\shortans{21}
\loigiai{
Ta có: $7 \times 3 = 21$.
}
\end{ex}

\begin{ex}
Số 8 được lấy 2 lần viết thành phép nhân là $8 \times 2$. Tính kết quả.
\shortans{16}
\loigiai{
Ta có: $8 \times 2 = 16$.
}
\end{ex}

\begin{ex}
Số 10 được lấy 4 lần viết thành phép nhân là $10 \times 4$. Tính kết quả.
\shortans{40}
\loigiai{
Ta có: $10 \times 4 = 40$.
}
\end{ex}

\begin{ex}
Mỗi đĩa có 3 quả cam, có 4 đĩa như thế. Hỏi có tất cả bao nhiêu quả cam?
\shortans{12}
\loigiai{
Có: $3 \times 4 = 12$ quả cam.
}
\end{ex}

\begin{ex}
Mỗi bạn có 2 chiếc bút chì, 5 bạn có bao nhiêu chiếc bút chì?
\shortans{10}
\loigiai{
Có: $2 \times 5 = 10$ chiếc bút.
}
\end{ex}

\begin{ex}
Mỗi con gà có 2 cái chân. Hỏi 6 con gà có tất cả bao nhiêu cái chân?
\shortans{12}
\loigiai{
Có: $2 \times 6 = 12$ cái chân.
}
\end{ex}

\begin{ex}
Mỗi con bò có 4 cái chân. Hỏi 5 con bò có tất cả bao nhiêu cái chân?
\shortans{20}
\loigiai{
Có: $4 \times 5 = 20$ cái chân.
}
\end{ex}

\begin{ex}
Mỗi tuần lễ bạn Nam đi học 5 ngày. Hỏi trong 3 tuần lễ bạn Nam đi học bao nhiêu ngày?
\shortans{15}
\loigiai{
Đi học: $5 \times 3 = 15$ ngày.
}
\end{ex}

\begin{ex}
Mỗi gói bánh có 5 chiếc bánh. Hỏi 4 gói bánh như thế có bao nhiêu chiếc bánh?
\shortans{20}
\loigiai{
Có: $5 \times 4 = 20$ chiếc bánh.
}
\end{ex}

\begin{ex}
Một bàn học có 2 học sinh ngồi. Hỏi 8 bàn học như thế có bao nhiêu học sinh ngồi?
\shortans{16}
\loigiai{
Có: $2 \times 8 = 16$ học sinh.
}
\end{ex}

\begin{ex}
Mỗi hàng cây có 6 cây. Hỏi 3 hàng cây như thế có tất cả bao nhiêu cây?
\shortans{18}
\loigiai{
Có: $6 \times 3 = 18$ cây.
}
\end{ex}

\begin{ex}
Điền số thích hợp: $2 + 2 + 2 + 2 + 2 = 2 \times \dots$
\shortans{5}
\loigiai{
Số 2 được lấy 5 lần nên số cần điền là 5.
}
\end{ex}

\begin{ex}
Điền số thích hợp: $5 + 5 + 5 + 5 = 5 \times \dots$
\shortans{4}
\loigiai{
Số 5 được lấy 4 lần nên số cần điền là 4.
}
\end{ex}

\begin{ex}
Điền số thích hợp: $3 \times 4 = 3 + 3 + 3 + \dots$
\shortans{3}
\loigiai{
Số cần điền là 3.
}
\end{ex}

\begin{ex}
Tính tổng $4 + 4 + 4 + 4 + 4$.
\shortans{20}
\loigiai{
Ta có $4 \times 5 = 20$.
}
\end{ex}

\begin{ex}
Tính tổng $10 + 10 + 10 + 10 + 10$.
\shortans{50}
\loigiai{
Ta có $10 \times 5 = 50$.
}
\end{ex}

\begin{ex}
Một chiếc xe ô tô con có 4 bánh xe. Hỏi 6 chiếc xe ô tô con như thế có bao nhiêu bánh xe?
\shortans{24}
\loigiai{
Có: $4 \times 6 = 24$ bánh xe.
}
\end{ex}

\begin{ex}
Mỗi lọ hoa cắm 5 bông hoa hồng. Hỏi 6 lọ hoa như thế cắm tất cả bao nhiêu bông hoa hồng?
\shortans{30}
\loigiai{
Có: $5 \times 6 = 30$ bông hoa.
}
\end{ex}

\begin{ex}
Số 2 được lấy 9 lần có giá trị bằng bao nhiêu?
\shortans{18}
\loigiai{
Giá trị: $2 \times 9 = 18$.
}
\end{ex}

\begin{ex}
Số 5 được lấy 7 lần có giá trị bằng bao nhiêu?
\shortans{35}
\loigiai{
Giá trị: $5 \times 7 = 35$.
}
\end{ex}
